% SPDX-License-Identifier: Apache-2.0
\section{\code{txhdl::regmap}}
\label{sec:r-regmap}

What a register map declared with \code{regmap!} stands for at run
time. The declaration is the macro: a peripheral says its map once, a
word index, a name, an access and a sentence a register, and under a
register its fields, each a name, a low bit, a width, an access, a
reset value and a sentence. An access is one of five: read and
written, read only, write only, write one to clear, or taken by the
read, the last for a register whose next read gives the next value,
as a receive queue's head does.

From that the macro writes four things. A module of constants a
program addresses by: the byte offset of each register and a
\code{Field} for each field. The read mux and the write enables a
unit's \code{run} calls, and the read enables beside them when the
declaration names a third function, which a register the read takes
needs. A function a field that takes it out of a written word, and
one a register that packs its fields into the word. And a table.
The lowering reads the same declaration in the file and inlines the
same functions into the netlist, so the map a program sees
and the map the hardware decodes are one text; the examples
document~\cite{examples} shows a peripheral written that way.
That holds for the decode and not for a field's reset, which the
header states and only the unit can hold: the unit builds the
register with it, and \code{Reg::reset\_mismatches} names every field
of a word read after reset that differs from it (issue 887).

This module is the data. \code{Field} is the constant a program uses:
its shift, its width and its reset, with \code{mask}, \code{get},
\code{set} and \code{with}, all of them \code{const}, so a field's
mask is a constant where a program wants one. \code{RegMap} is the
table, its \code{Reg}s and their \code{FieldInfo}s, which
\code{//tools/regmap} writes as a C header, a define a register with
its offset and four a field, as a device tree node at a base, or as a
Rust module of the same constants, and \code{//tools/datasheet}
writes as a sheet's register table, a row a register and a row a
field under it.

The software takes its offsets from those outputs rather than typing
them (issue 709). \code{regmap all rs} writes every map into one file,
and the crate \code{vreteno\_regs} under
\code{//cpu/vreteno/rust/regs} is that file: it depends on nothing, so
the core's firmware, which has no standard library and cannot depend
on the parts, names it and addresses a register as a base plus
\code{vreteno\_regs::uart::STATUS}. The Zephyr drivers include C
headers written the same way, committed under
\code{zephyr/include/vreteno/regs}, since a \code{west build} of the
port runs no Bazel; \code{//zephyr:regs\_test} holds each committed
header to its map, and \code{bazel run //zephyr:regs\_update} writes
them again after a map changes. A register moved in the hardware then
moves in the software, or a test says which header is stale.

\lstinputlisting[caption={\code{lib/src/regmap.rs}.},label={lst:r-regmap}]{lib/src/regmap.rs}
